\section{Introduction}
\label{sec::intro}
The advent of artificial intelligence (AI) has indeed transformed the research paradigm in traditional economics~\cite{jorgenson2001information}. In the digital economy era, individual economic behavior can be elaborately recorded and analyzed, leading to a growing emphasis on data-driven modeling in macroeconomics. These data sources may include online consumption, social media activity, and more, allowing economists to better understand individual economic behavior rather than relying solely on traditional macroeconomic indicators~\cite{schorfheide2015real}.
Research in macroeconomics aims to analyze and predict economic variables quantitatively. Early empirical statistical models, such as the Phelps Model~\cite{phelps1967phillips}, and the work of Kydland and Prescott~\cite{kydland1982time}, focused on data-driven analysis and policy outcome prediction but struggled to handle significant shocks. Later, Dynamic Stochastic General Equilibrium (DSGE) models~\cite{christiano2005nominal} were introduced to capture economic dynamics and shocks but assumed a perfect world.

In the last two decades, agent-based modeling (ABM) has emerged as a promising paradigm for simulating macroeconomics from the bottom up, allowing diverse agents to interact without assuming a predetermined equilibrium~\cite{farmer2009economy}. The evolution of ABM in macroeconomics can be divided into two stages. Early models~\cite{tesfatsion2006handbook,brock1998heterogeneous} relied on predetermined rules but made oversimplified assumptions about agent behaviors. Later learning-based models~\cite{trott2021building,zheng2022ai,mi2023taxai}, were trained with large-scale behavioral data. Essentially, these models transform multifaceted economic factors into agent decisions.
However, customizing decision-making mechanisms for each agent presents substantial difficulties. Using customized rules requires extensive expert knowledge and sophisticated model calibration~\cite{windrum2007empirical}. Similarly, employing customized neural networks leads to a sharp increase in network parameters and difficulties in model training~\cite{mi2023taxai}. These difficulties raise the challenge of modeling agent heterogeneity. Additionally, existing models typically focus on the individual situations of the current period, making it difficult to consider past periods, their variations, and multifaceted macroeconomic factors. This complicates the modeling of the influence of macroeconomic trends and dynamics.

Recently, the field of AI has witnessed the rise of large language models (LLMs)~\cite{zhao2023survey}. Leveraging this advancement, LLM-empowered agents have demonstrated capabilities in reasoning, planning, and decision-making. In this work, we design \textbf{EconAgent}, a LLM-empowered agent with human-like characteristics for macroeconomic simulations. We first construct a simulation environment that includes labor and consumption market dynamics driven by agents' decisions on working and consumption, as well as fiscal and monetary policies. Using a perception module that targets agent profiles and mirrors real-world economic situations, we create heterogeneous agents automatically exhibiting different decision-making mechanisms. Additionally, we model the influence of macroeconomic trends with a memory module, enabling agents to reflect on past individual experiences and market dynamics.

In our experiments, classic macroeconomic phenomena, such as inflation in the consumption market and the unemployment rate in the labor market, are reproduced more reasonably compared to traditional rule-based or learning-based agents. We also observe that the EconAgent exhibits human-like decision-making patterns, reflecting the complexities of human economic behavior. These agents demonstrate swift adaptability in response to changes in the internal and external environment.

The contributions of this work can be summarized as follows:
\begin{itemize}[leftmargin=*]
    \item We take the first step to integrate LLMs into the domain of macroeconomic simulations, bridging two seemingly disparate fields into a cohesive research avenue.
    \item We conduct macroeconomic simulations in our constructed environment driven by EconAgent, resulting in the emergence of classic macroeconomic phenomena and regularities.
    \item The results show that our approach not only enhances the realism and depth of macroeconomic simulations but also provides a promising avenue for future research, potentially reshaping how we study and understand the intricacies of macroeconomics.
\end{itemize}